\documentclass[10pt,a4paper]{article}

\usepackage[utf8]{inputenc}
\usepackage[T1]{fontenc}
\usepackage[brazil]{babel}
\usepackage[margin=1.8cm]{geometry}
\usepackage{xcolor}
\usepackage{booktabs}
\usepackage{enumitem}
\usepackage[hidelinks]{hyperref}

\setlist{noitemsep,topsep=2pt,leftmargin=1.4em}
\setlength{\parskip}{2pt}
\setlength{\parindent}{0pt}

\title{\vspace{-1.0cm}\textbf{Java Event Planner} \\[1mm]
\large Relatório do Projeto Final}
\author{}
\date{}

\begin{document}
\maketitle
\vspace{-1.4cm}

% ---------------------------------------------------------------- group id
\begin{center}
\textbf{SCC0204 — Programação Orientada a Objetos} \quad ICMC — USP \\[1mm]
\begin{tabular}{ll}
\toprule
\textbf{Integrante} & \textbf{Nº USP} \\
\midrule
Fabio Kauê Araujo da Silva & 16311045 \\
Kainã Alves Tureso & 15466391 \\
Luís Henrique de Queiroz Veras & 14592414 \\
\bottomrule
\end{tabular} \\[1mm]
Repositório: \url{https://github.com/kainaas/Trabalho_POO}
\end{center}

% ---------------------------------------------------------------- 1
\section{Requisitos}
O \textbf{Java Event Planner} é um aplicativo \emph{desktop} (Java + Swing) para
organização de compromissos em um calendário. \textbf{Requisitos funcionais:}
\begin{itemize}
  \item Cadastro, edição e exclusão de eventos com título, data, hora, duração,
        local, descrição e categoria (reunião, aniversário, compromisso ou outro).
  \item Eventos recorrentes (diário, semanal ou mensal), com a opção de
        editar/excluir \emph{somente uma ocorrência} ou \emph{a série inteira}.
  \item Participantes (nome e e-mail) e lembretes (10\,min, 1\,h ou 1\,dia antes),
        exibidos em tempo de execução.
  \item Busca por palavra-chave (título, descrição, local ou categoria).
  \item Persistência dos eventos em arquivo, recarregados na abertura do programa.
\end{itemize}
\textbf{Requisitos adicionais} levantados pela nossa implementação:
\begin{itemize}
  \item \textbf{Quatro modos de visualização} — Dia (agenda por hora), Semana (sete
        colunas), Mês (grade) e Ano (os doze meses) — trocados na barra superior.
  \item \textbf{Tema claro/escuro} aplicado a toda a interface, com ícones próprios
        para cada tema e botões totalmente repintados (sem ``faixas brancas'').
  \item \textbf{Indicador da quantidade de eventos} de cada dia, na própria célula.
  \item \textbf{Seletor de data por calendário} no formulário de evento.
  \item \textbf{Robustez na leitura}: linhas inválidas do arquivo são ignoradas.
\end{itemize}

% ---------------------------------------------------------------- 2
\section{Descrição do Projeto}
A aplicação segue o padrão \textbf{MVC}, refletido na divisão em pacotes:
\textbf{Model} (dados e regras: \texttt{CalendarModel}, \texttt{Event},
\texttt{Attendee}, \texttt{Recurrence}, \texttt{EventStorage}, \texttt{AbstractModel}
e as exceções), \textbf{View} (telas Swing: \texttt{CalendarView},
\texttt{TopPanel}, os quatro modos \texttt{DayViewPanel}/\texttt{WeekViewPanel}/%
\texttt{CalendarMonthPanel}/\texttt{YearViewPanel}, \texttt{DayEventsPanel},
\texttt{EventFormDialog}, \texttt{DatePickerDialog}, \texttt{ReminderService},
\texttt{Icons}, \texttt{ThemeColors}, \texttt{ViewMode}), \textbf{Controller}
(\texttt{Controller}) e \textbf{Main}.

\textbf{Observer.} \texttt{CalendarModel} herda de \texttt{AbstractModel}, que
encapsula um \texttt{PropertyChangeSupport}. Toda alteração relevante (novo evento,
troca de data, de modo de visualização ou de tema) dispara uma notificação;
a \texttt{CalendarView} implementa \texttt{PropertyChangeListener} e redesenha os
painéis, sem que o modelo conheça a interface.

\textbf{Visualizações.} O centro da janela troca o painel conforme o \texttt{ViewMode}
atual (interface \texttt{CalendarSubView}). \textbf{Recorrência:} guarda-se apenas o
evento ``mestre'' e o método \texttt{occursOn(dia)} calcula sob demanda se ele cai em
uma data, usando um conjunto de \emph{datas de exceção} para remoções pontuais. O
diagrama de classes está em \texttt{docs/uml.puml} e a documentação gerada
(\texttt{javadoc}) em \texttt{docs/javadoc}.

% ---------------------------------------------------------------- 3
\section{Comentários sobre o Código}
O código está documentado no padrão \textbf{Javadoc} (em inglês) nas classes e
métodos não triviais. Conceitos de POO aplicados:
\begin{itemize}
  \item \textbf{Encapsulamento:} campos \texttt{private} expostos por
        \emph{getters/setters}; a validação de um \texttt{Event} é centralizada no
        \texttt{Controller}.
  \item \textbf{Herança:} \texttt{CalendarModel} de \texttt{AbstractModel}; as
        exceções de \texttt{Exception}; os painéis de \texttt{JPanel}, a janela de
        \texttt{JFrame} e os formulários de \texttt{JDialog}.
  \item \textbf{Polimorfismo:} \texttt{toString()} sobrescrito em \texttt{Recurrence},
        \texttt{Attendee} e \texttt{ViewMode}; \texttt{CalendarView} implementa
        \texttt{PropertyChangeListener}; \texttt{ReminderService} implementa
        \texttt{Runnable}; os quatro modos implementam \texttt{CalendarSubView}.
  \item \textbf{Composição:} um \texttt{CalendarModel} é dono da lista de
        \texttt{Event}, e cada \texttt{Event} é dono de seus \texttt{Attendee}.
\end{itemize}
As cores de cada tema ficam em \texttt{ThemeColors}, cujo método \texttt{styleButton}
repinta os botões de forma consistente mesmo no \emph{Look and Feel} padrão (Metal),
evitando as ``faixas brancas'' no modo escuro.

% ---------------------------------------------------------------- 4
\section{Plano de Testes}
A lógica de negócio é coberta por uma suíte automatizada em \textbf{JUnit~5}
(\texttt{test/}, 28 casos), executável com \texttt{test/run-tests.cmd}:
\begin{itemize}
  \item \texttt{EventTest} — \texttt{occursOn} para os quatro tipos de recorrência,
        datas de exceção e cálculo do fim do evento.
  \item \texttt{CalendarModelTest} — título único, contagem por dia, ordenação
        cronológica, busca e remoção de ocorrência.
  \item \texttt{ControllerTest} — validações (título vazio, duração $\le 0$, data
        inválida), bloqueio de duplicata e gravação em disco.
  \item \texttt{EventStorageTest} — ida-e-volta de gravação/leitura, tolerância a
        linhas corrompidas e arquivo inexistente.
  \item \texttt{ModelValueTest} — \texttt{toString} de \texttt{Attendee} e
        \texttt{Recurrence}.
\end{itemize}
Complementarmente, foram feitos testes manuais de interface: troca de tema, troca de
modo de visualização, lembretes e o seletor de data por calendário.

% ---------------------------------------------------------------- 5
\section{Resultados dos Testes}
Todos os 28 testes passam (Java~25, JUnit~5.8.2). Saída resumida do console:
\begin{center}\small
\texttt{[\ 28 tests found\ ] [\ 28 tests started\ ] [\ 28 tests successful\ ] [\ 0 tests failed\ ]}
\end{center}
Os testes manuais de interface também passaram: o tema é aplicado a botões, ícones,
células e campos (sem ``faixas brancas''); os quatro modos de visualização funcionam;
a persistência é estável e um arquivo corrompido não derruba o programa.

% ---------------------------------------------------------------- 6
\section{Procedimentos de Build}
Não há bibliotecas externas — apenas o JDK (testado no Java~25; JDK~17+ deve funcionar).
\begin{enumerate}
  \item \textbf{Instalar o JDK} (ex.: Eclipse Temurin, \url{https://adoptium.net});
        verifique com \texttt{java -version} e \texttt{javac -version}.
  \item \textbf{Obter o código:} \texttt{git clone https://github.com/kainaas/Trabalho\_POO.git}
        e \texttt{cd Trabalho\_POO/src}.
  \item \textbf{Compilar} (Windows: \texttt{compile.cmd}):
        \texttt{javac -encoding UTF-8 -d bin ./Controller/*.java ./Model/*.java ./View/*.java ./Main/*.java}
  \item \textbf{Executar} (Windows: \texttt{Run.cmd}): \texttt{cd bin \&\& java Main.Main}.
        O diretório de trabalho deve ser \texttt{bin} para resolver os ícones em
        \texttt{bin/iconImages}; \texttt{events.txt} é criado ali no primeiro salvamento.
  \item \textbf{Testes (opcional):} com \texttt{junit-platform-console-standalone-1.8.2.jar}
        na raiz, rode \texttt{test/run-tests.cmd}.
\end{enumerate}

% ---------------------------------------------------------------- 7
\section{Problemas}
\begin{itemize}
  \item \textbf{Recorrência.} A primeira ideia (uma cópia por dia) gastava memória;
        a solução foi o evento ``mestre'' com \texttt{occursOn} sob demanda e um
        conjunto de datas de exceção.
  \item \textbf{Título único vs.\ editar uma ocorrência.} A regra de título único
        conflitava com a edição de ``somente esta'' ocorrência (que gera um avulso
        de mesmo título); resolvido com \texttt{addEventForced}, usado só nesse caso.
  \item \textbf{Formato do arquivo.} Persistir listas exigiu separadores em dois
        níveis (TAB e vírgula) e leitura tolerante a linhas inválidas.
  \item \textbf{Concorrência.} O \texttt{ReminderService} roda em uma \emph{thread}
        \emph{daemon} e os avisos usam \texttt{SwingUtilities.invokeLater}.
\end{itemize}

% ---------------------------------------------------------------- 8
\section{Comentários}
O projeto foi desenvolvido de forma incremental e versionado no GitHub. A arquitetura
MVC com Observer manteve a interface desacoplada das regras de negócio, o que facilitou
adicionar recursos (tema escuro, indicador de eventos, modos de visualização e seletor
de data) sem reescrever as telas. Conforme o \emph{peer review} da disciplina, os
apontamentos do grupo revisor serão documentados e respondidos nesta seção na versão
final.

\end{document}
